\documentclass[../main.tex]{subfiles}
\begin{document}

\chapter{Fault Tolerance}
\lessoninfo{
  video={Lesson 17, videos 1--32},
  textbook={H\&P \cite{hennessy2017} §1.7, §D.2--D.3},
  keywords={dependability, fault, error, failure, reliability, availability, MTTF, MTTR, fault tolerance, checkpointing, N-module redundancy, parity, ECC, RAID},
}

Lesson 16 described the devices that store data persistently.  Every
component of a computer, storage devices included, eventually deviates from
its intended behavior, and a system that returns wrong results or loses data
whenever one component does so cannot be relied on.  This lesson introduces
the vocabulary for reasoning about such problems---faults, errors and
failures; reliability and availability---classifies the faults that occur in
practice, and presents the general techniques for tolerating them.  The
second half applies these ideas to disks: redundant arrays of independent
disks (RAID) combine several disks into one, and their redundant
organizations keep data intact and accessible even when an entire disk
fails.

\section{Dependability}
\label{sec:l17-dependability}

A computer is useful only if its users can trust the service it
provides.\source{Lesson 17, videos 1--2; H\&P §1.7;
\cite{avizienis2004}.}  To make this precise, two descriptions of the
system's behavior are distinguished.  The \emph{specified service} is the
behavior the system should have according to its specification; the
\emph{delivered service} is the behavior it actually exhibits.

\begin{definition}[Dependability]
  The \term{dependability}[ディペンダビリティ] of a system is the quality of
  its delivered service that justifies relying on the system to provide that
  service.
\end{definition}

A system is built from components, or \emph{modules}: processors, memory
modules, disks and, at a finer level, the circuits and program functions
inside them.  Each module has a specified behavior of its own.
Dependability problems begin when some module deviates from its
specification; the next section describes how such a deviation can, but need
not, turn into a deviation of the whole system.

\section{Faults, errors and failures}
\label{sec:l17-fault-error-failure}

A problem that starts inside one module passes through up to three stages
before the users of the system can observe it.\source{Lesson 17, videos
3--5; H\&P §1.7, §D.3.}  The stages are named separately because a
dependable system tries to stop the problem at one of the transitions
between them.

\begin{definition}[Fault, error, failure]
  \label{def:l17-fault-error-failure}
  A \term{fault}[フォールト] is a deviation of a module from its specified
  behavior, such as a defective transistor or a mistake in a program.

  An \term{error}[エラー] exists when the actual behavior inside the
  system---its internal state---differs from the specified behavior, such as
  a wrong value stored in a variable.

  A \term{failure}[故障] occurs when the delivered service of the system
  deviates from the specified service.
\end{definition}

A fault can exist for a long time without any effect.  Until the faulty
module is used in a way that exposes the deviation, the fault is a
\emph{latent} error.  When the module is used in that way, the fault is
\emph{activated} and becomes an \emph{effective} error: the internal state
is now wrong.  An effective error causes a failure only if it propagates to
the service the system delivers.  Not every fault is activated, and not every error becomes a
failure (\cref{fig:l17-fault-chain}).

\begin{example}
  \label{ex:l17-fault-chain}
  A function that multiplies two integers contains a programming mistake:
  it returns the correct product for all arguments except 6 and 9, for
  which it returns 45 instead of 54.  The mistake is a fault from the moment
  the function is written, and it stays latent as long as no caller
  multiplies 6 by 9.  When a program calls the function with these arguments
  and stores the result, the fault has been activated and the program state
  contains an error.  If the result is the number of seats a booking system
  reserves, 45 seats are reserved instead of 54: the error has reached the
  delivered service, and a failure has occurred.  If instead the program
  only tests whether the product is positive, the wrong value leads to the
  same decision as the correct one, and no failure occurs.
\end{example}

\begin{figure}[htbp]
  \centering
  % Fault -> error -> failure.  An error that leaves a module is a failure of
% that module and, at the same time, a fault of the enclosing system; an error
% that reaches the system's delivered service is a failure of the system.
\begin{tikzpicture}[x=1cm,y=1cm, font=\footnotesize\sffamily, >={Stealth[length=1.8mm]},
  st/.style={draw, rounded corners=2pt, fill=white, minimum height=0.65cm,
    minimum width=1.2cm, inner sep=2pt},
  lab/.style={font=\scriptsize, text=ln-muted, align=center}]
  % boundaries
  \draw[ln-rule, fill=ln-box] (0,0) rectangle (9.0,3.2);
  \node[anchor=north west, font=\scriptsize\sffamily\bfseries, text=ln-accent]
    at (0.05,3.15) {\iflangja{システム}{System}};
  \draw[ln-rule, fill=ln-accent!10] (0.25,1.05) rectangle (4.8,2.75);
  \node[anchor=north west, font=\scriptsize\sffamily\bfseries, text=ln-accent]
    at (0.3,2.7) {\iflangja{モジュール}{Module}};
  % chain
  \node[st] (f)  at (1.1,1.75) {\iflangja{フォールト}{fault}};
  \node[st] (e1) at (3.85,1.75) {\iflangja{エラー}{error}};
  \node[st] (e2) at (7.0,1.75) {\iflangja{エラー}{error}};
  \node[st, draw=ln-caution, text=ln-caution] (fl) at (10.6,1.75) {\iflangja{故障}{failure}};
  \draw[->] (f) -- node[lab, above=1pt] {\iflangja{活性化}{activated}} (e1);
  \draw[->] (e1) -- node[lab, above=1pt, pos=0.55] {\iflangja{伝播}{propagates}} (e2);
  \draw[->] (e2) -- node[lab, above=1pt, pos=0.28] {\iflangja{伝播}{propagates}} (fl);
  % boundary crossings
  \fill[ln-accent] (4.8,1.75) circle (1.6pt);
  \fill[ln-caution] (9.0,1.75) circle (1.6pt);
  \node[lab, anchor=north] at (4.8,0.95)
    {\iflangja{モジュールの故障＝\\システムのフォールト}{module failure =\\fault of the system}};
  \node[lab, anchor=north] at (10.6,1.3)
    {\iflangja{提供サービスが\\仕様から逸脱}{delivered service\\deviates from\\specified service}};
  \node[lab, anchor=north] at (1.1,1.38) {\iflangja{（潜在）}{(latent)}};
\end{tikzpicture}

  \caption{A fault becomes an error when activated.  An error that leaves a
    module is a failure of the module and, at the same time, a fault of the
    system that contains it; an error that reaches the delivered service is a
    failure of the system.}
  \label{fig:l17-fault-chain}
\end{figure}

The three terms are relative to the boundary under consideration.  A failure
of a module is, seen from the system that contains the module, a fault of
that system.  A disk whose head has crashed has failed as a disk; for an
array of disks that stores redundant copies of the data (\cref{sec:l17-raid}),
the same event is a fault that need not lead to a failure of the array.

\section{Reliability and availability}
\label{sec:l17-reliability}

With respect to its specified service, a system is at any time in one of
two states.\source{Lesson 17, videos 6--7; H\&P §1.7.}  In \emph{service
accomplishment} the delivered service matches the specified service; in
\emph{service interruption} it does not.  A failure moves the system from
the first state to the second, and a restoration (repair) moves it back
(\cref{fig:l17-service-timeline}).  The time spent in each state gives two
different measures of dependability.

\begin{figure}[htbp]
  \centering
  % Two service states and a timeline alternating between them.  Up-times are
% times to failure, down-times are times to repair.
\begin{tikzpicture}[x=1cm,y=1cm, font=\footnotesize\sffamily, >={Stealth[length=1.8mm]},
  st/.style={draw, rounded corners=3pt, minimum height=0.65cm, inner sep=3pt},
  lab/.style={font=\scriptsize, text=ln-muted, align=center}]
  % states
  \node[st, fill=ln-accent!15] (up) at (2.6,2.3) {\iflangja{サービス遂行}{service accomplishment}};
  \node[st, fill=ln-caution!15] (down) at (8.4,2.3) {\iflangja{サービス中断}{service interruption}};
  \draw[->] (up) to[bend left=12] node[lab, above] {\iflangja{故障}{failure}} (down);
  \draw[->] (down) to[bend left=12] node[lab, below] {\iflangja{修復}{restoration}} (up);
  % timeline
  \foreach \a/\b/\n in {0/3.1/1, 3.6/6.4/2, 7.2/10.6/3} {
    \fill[ln-accent!15] (\a,0) rectangle (\b,0.45);
    \draw[|<->|, ln-muted] (\a,0.7) -- node[above, font=\scriptsize] {$U_\n$} (\b,0.7);
  }
  \foreach \a/\b/\n in {3.1/3.6/1, 6.4/7.2/2, 10.6/11.1/3} {
    \fill[ln-caution!40] (\a,0) rectangle (\b,0.45);
    \node[font=\scriptsize, below] at ({(\a+\b)/2},-0.05) {$D_\n$};
  }
  \draw (0,0) rectangle (11.1,0.45);
  \draw[->] (0,-0.55) -- (11.4,-0.55) node[pos=1, below left, lab] {\iflangja{時間}{time}};
  \node[lab, anchor=north] at (5.55,-0.75)
    {MTTF $=$ \iflangja{$U_i$ の平均}{mean of $U_i$}\qquad
     MTTR $=$ \iflangja{$D_i$ の平均}{mean of $D_i$}};
\end{tikzpicture}

  \caption{The two service states and a history that alternates between
    them.  $U_i$ are the periods of service accomplishment, $D_i$ the periods
    of service interruption.}
  \label{fig:l17-service-timeline}
\end{figure}

\begin{definition}[Reliability]
  The \term{reliability}[信頼性] of a system is a measure of
  \emph{continuous} service accomplishment.  It is usually quantified by the
  \term{mean time to failure}[平均故障時間]\index{MTTF|see{mean time to failure}}
  (MTTF): the average length of a period of service accomplishment, i.e.\
  the average time from the start of service until the next failure.
\end{definition}

\begin{definition}[Availability]
  The \term{availability}[可用性] of a system is the fraction of time it
  spends in service accomplishment.  With the \term{mean time to
  repair}[平均修復時間]\index{MTTR|see{mean time to repair}} (MTTR), the
  average length of a period of service interruption,
  \begin{equation}
    \text{availability} \;=\; \frac{\mathrm{MTTF}}{\mathrm{MTTF}+\mathrm{MTTR}} .
    \label{eq:l17-availability}
  \end{equation}
\end{definition}

The sum of the two means is the average length of one complete
failure--repair cycle, the \term{mean time between failures}[平均故障間隔]%
\index{MTBF|see{mean time between failures}}:
\begin{equation}
  \mathrm{MTBF} \;=\; \mathrm{MTTF} + \mathrm{MTTR} .
  \label{eq:l17-mtbf}
\end{equation}
When repairs are short compared with the periods of service, the MTBF and
the MTTF are nearly equal.

\begin{example}
  \label{ex:l17-availability}
  A server runs for \SI{1200}{\hour}, is down for \SI{6}{\hour}, runs for
  \SI{900}{\hour}, is down for \SI{10}{\hour}, runs for \SI{1500}{\hour} and is
  down for \SI{8}{\hour}.  Its MTTF is $(1200+900+1500)/3 = \SI{1200}{\hour}$,
  its MTTR is $(6+10+8)/3 = \SI{8}{\hour}$, its MTBF is \SI{1208}{\hour}, and
  its availability is $1200/1208 \approx \SI{99.34}{\percent}$.  A second
  server that fails ten times as often (MTTF $=\SI{120}{\hour}$) but is
  restored within \SI{0.8}{\hour} on average has the same availability but
  only one tenth of the MTTF, so it is far less reliable.
\end{example}

\section{Kinds of faults}
\label{sec:l17-kinds}

Faults are classified in two ways: by their cause and by how long they
last.\source{Lesson 17, videos 8--9; H\&P §D.3.}

\subsection{By cause}

\begin{description}
  \item[Hardware] A \term{hardware fault}[ハードウェアフォールト] occurs when
    hardware fails to perform as designed, e.g.\ a transistor that no longer
    switches or a disk head that touches the platter.
  \item[Design] A \term{design fault}[設計フォールト] is a mistake in the
    design itself: a software bug, or a hardware design error such as the
    division bug of the early Intel Pentium, which was present in every chip
    built from the flawed design.
  \item[Operation] An \term{operation fault}[運用フォールト] is a mistake
    made by the operators or users of the system, e.g.\ deleting the wrong
    file or disconnecting the wrong cable.
  \item[Environment] An \term{environmental fault}[環境フォールト] originates
    outside the system: fire, flooding, a power outage, extreme temperature,
    or deliberate sabotage.
\end{description}

\subsection{By duration}

\begin{description}
  \item[Permanent] A \term{permanent fault}[永久フォールト] persists once it
    has occurred until the faulty part is repaired or replaced, e.g.\ a
    chip damaged by overheating.
  \item[Intermittent] An \term{intermittent fault}[間欠フォールト] lasts for
    a while, disappears, and recurs, e.g.\ a loose connector, or a processor
    running above its rated clock frequency that miscalculates only when it
    becomes hot.
  \item[Transient] A \term{transient fault}[過渡フォールト] occurs once and
    goes away by itself, e.g.\ a bit of memory flipped by a high-energy
    particle; rewriting the bit restores correct operation.
\end{description}

The two classifications describe different properties of a fault.  A
hardware fault, for instance, can be permanent (a burnt-out chip),
intermittent (a cracked solder joint that loses contact only when warm) or
transient (a flipped memory bit).  A design fault, in contrast, is present in
every copy of the design until the design is corrected, so it is permanent,
even if it is activated only occasionally.

\section{Improving reliability and availability}
\label{sec:l17-improving}

There are three ways to make a system more dependable.\source{Lesson 17,
video 10; H\&P §1.7.}

\begin{description}
  \item[Avoid faults] Preventing faults from occurring in the first place is
    called \term{fault avoidance}[フォールト回避]: for example, using
    components with lower failure rates, keeping them cool, or keeping
    liquids away from the machines.
  \item[Tolerate faults] Preventing faults from becoming failures is called
    \term{fault tolerance}[フォールトトレランス]: the system keeps
    delivering its specified service although some of its modules are
    faulty.  Fault tolerance is built on redundancy, for instance
    error-correcting codes in memory (\cref{sec:l17-codes}) or redundant
    disks (\cref{sec:l17-raid}).
  \item[Repair faster] Shortening the time from a failure to the restoration
    of service, e.g.\ by keeping spare parts on site.
\end{description}

\begin{keypoint}
  Fault avoidance and fault tolerance reduce the number of failures, so they
  raise the MTTF and thereby both reliability and availability.  Faster
  repair lowers only the MTTR; for a system that stops whenever a module
  fails it improves availability but not reliability.
\end{keypoint}

\section{Fault-tolerance techniques}
\label{sec:l17-techniques}

General-purpose fault tolerance combines two ingredients: a way to
\emph{detect} errors and a way to \emph{recover} from them.\source{Lesson
17, videos 11--13.}

\subsection{Checkpointing}

Saving the complete state of a system at regular intervals, so that
execution can later resume from a saved state, is called
\term{checkpointing}[チェックポインティング].  When an error is detected,
the system restores the most recent checkpoint and executes again from
there.  The re-execution is likely to succeed if the fault was transient,
and often if it was intermittent; a permanent fault is activated again and
cannot be tolerated this way.  Recovery also takes time.  If restoring the
checkpoint and repeating the lost work is quick, users perceive only a delay;
if it takes too long, the recovery itself counts as a service interruption.

\subsection{Redundancy}

\begin{definition}[$N$-module redundancy]
  In \term[N-module redundancy]{$N$-module redundancy}[N モジュール冗長]
  $N$ identical modules perform the same computation on the same inputs, and
  their outputs are compared or put to a majority vote
  (\cref{fig:l17-nmr}).
\end{definition}

\begin{figure}[htbp]
  \centering
  % N-module redundancy: (a) two modules and a comparator detect an error;
% (b) three modules and a majority voter mask the error of one module.
\begin{tikzpicture}[x=1cm,y=1cm, font=\footnotesize\sffamily, >={Stealth[length=1.8mm]},
  mod/.style={draw, fill=ln-box, minimum width=1.1cm, minimum height=0.55cm, inner sep=2pt},
  lab/.style={font=\scriptsize, text=ln-muted, align=center}]
  % ---------------- (a) dual
  \node[anchor=west, font=\footnotesize\bfseries] at (-0.5,1.55) {(a) $N=2$};
  \draw (-0.5,0) node[left, lab] {\iflangja{入力}{in}} -- (0.3,0);
  \fill (0.3,0) circle (1.2pt);
  \node[mod] (a1) at (1.3,0.55) {M1};
  \node[mod] (a2) at (1.3,-0.55) {M2};
  \draw[->] (0.3,0) |- (a1.west);
  \draw[->] (0.3,0) |- (a2.west);
  \node[draw, circle, fill=ln-accent!15, minimum size=0.7cm, inner sep=0pt] (cmp) at (3.1,0) {$=$};
  \draw[->] (a1.east) -| (cmp.north);
  \draw[->] (a2.east) -| (cmp.south);
  \draw[->] (cmp.east) -- (4.0,0) node[right, lab] {\iflangja{出力}{out}};
  \node[lab, anchor=north] at (3.1,-0.9)
    {\iflangja{一致: 結果は有効\\不一致: エラー検出}{equal: result valid\\differ: error detected}};
  % ---------------- (b) triple
  \begin{scope}[xshift=6.0cm]
  \node[anchor=west, font=\footnotesize\bfseries] at (-0.5,1.55) {(b) $N=3$};
  \draw (-0.5,0) node[left, lab] {\iflangja{入力}{in}} -- (0.3,0);
  \fill (0.3,0) circle (1.2pt);
  \node[mod] (b1) at (1.3,0.8) {M1};
  \node[mod] (b2) at (1.3,0) {M2};
  \node[mod] (b3) at (1.3,-0.8) {M3};
  \draw[->] (0.3,0) |- (b1.west);
  \draw[->] (0.3,0) -- (b2.west);
  \draw[->] (0.3,0) |- (b3.west);
  \node[draw, fill=ln-accent!15, minimum width=0.9cm, minimum height=2.1cm, align=center, inner sep=1pt] (vote) at (3.0,0)
    {\iflangja{多数決}{vote}};
  \draw[->] (b1.east) -- (b1.east -| vote.west);
  \draw[->] (b2.east) -- (vote.west);
  \draw[->] (b3.east) -- (b3.east -| vote.west);
  \draw[->] (vote.east) -- (4.1,0) node[right, lab] {\iflangja{出力}{out}};
  \end{scope}
\end{tikzpicture}

  \caption{$N$-module redundancy.  (a) Two modules and a comparator detect
    that one module is faulty, but not which.  (b) Three modules and a
    majority voter mask the wrong output of any one module.}
  \label{fig:l17-nmr}
\end{figure}

With two modules (\emph{dual-module redundancy}), a mismatch reveals that
one of them is faulty, but not which one, so the error is detected but
cannot be corrected by the comparison alone.  Combined with checkpointing,
the system rolls back and repeats the computation.  With three modules
(\emph{triple-module redundancy}), the correct result is the one produced by
at least two, and a single faulty module is outvoted without any
rollback---even if its fault is permanent.  In general a majority vote among
$N$ modules masks up to $\lfloor (N-1)/2 \rfloor$ faulty modules
(\cref{tab:l17-nmr}).  The price is $N$ times the hardware and
energy.\caution{All $N$ modules share one design, so a design fault is not
tolerated.}

\begin{table}[htbp]
  \centering\small
  \caption{Faulty modules handled by $N$-module redundancy.}
  \label{tab:l17-nmr}
  \begin{tabular}{@{}cp{0.75\linewidth}@{}}
    \toprule
    $N$ & Effect of faulty modules \\
    \midrule
    2 & one faulty module is detected (the outputs differ) but not corrected \\
    3 & one faulty module is outvoted 2:1 \\
    5 & one faulty module is outvoted 4:1, two are outvoted 3:2 \\
    \bottomrule
  \end{tabular}
\end{table}

With one faulty module, a five-module system still masks one more fault.
With two faulty modules, the three correct ones still outvote them 3:2, but
no margin is left: a third faulty module would outvote the two correct ones,
and the vote would produce a wrong result.  Operation is therefore stopped
before that can happen.  The Space Shuttle, with its five flight computers,
followed this policy: after one computer failed the mission continued
normally, and after a second one failed it was cut short.

\section{Codes for memory and storage}
\label{sec:l17-codes}

For memory and storage, replicating whole modules is overkill.\source{Lesson
17, video 14; H\&P ch.~2.}  A fault typically corrupts a few of the
stored bits, and those can be protected with far fewer extra bits by an
error detection or correction code.

\begin{definition}[Parity]
  The \term{parity}[パリティ] of a group of data bits $d_1,\dots,d_n$ is
  their exclusive OR,
  \begin{equation}
    p \;=\; d_1 \oplus d_2 \oplus \dots \oplus d_n ,
    \label{eq:l17-parity}
  \end{equation}
  and a parity bit is one extra bit that stores $p$ together with the data.
\end{definition}

When the data is read, the parity is computed again and compared with the
stored bit.  Flipping any single bit, including the parity bit itself,
changes the result, so every single-bit error is detected.  The parity does
not reveal which bit flipped, so the error cannot be corrected, and two
flipped bits cancel each other and go unnoticed.  For example, the byte
\texttt{10110010} has four ones and parity 0; if it is read back as
\texttt{10100010}, the recomputed parity is 1 and the error is detected.

\begin{definition}[Error-correcting code]
  An \term{error-correcting code}[誤り訂正符号]\index{ECC|see{error-correcting code}}
  (ECC) adds enough check bits to a group of data bits that a limited
  number of flipped bits can be both located and corrected.
\end{definition}

Memory modules with ECC typically use a SECDED code (single error
correction, double error detection), for example 8 check bits for every 64
data bits: any single flipped bit in a word is corrected, and any two
flipped bits are detected.  Disks protect each sector with considerably
stronger codes, such as Reed--Solomon codes, that correct many flipped bits,
in particular bursts of adjacent bits caused by defects on the recording
surface.  These codes handle damaged bits within a sector; they cannot help
when a whole disk stops working.  That case calls for redundancy across
disks.

\section{Redundant arrays of independent disks}
\label{sec:l17-raid}

A \term{redundant array of independent
disks}[RAID]\index{RAID|see{redundant array of independent disks}} (RAID) is
a set of disks that together play the role of a single disk, one that can be
larger, faster or more reliable than a single disk.\source{Lesson 17, video
15; H\&P §D.2; \cite{patterson1988}.}  An array has two goals:
\begin{itemize}
  \item better performance than a single disk, and
  \item normal read and write service even when a sector goes bad or an
    entire disk fails.
\end{itemize}
Not every organization pursues both goals: RAID 0 (\cref{sec:l17-raid0})
aims only at performance.

Each disk in the array still protects its sectors with its own code
(\cref{sec:l17-codes}).  When a sector cannot be read correctly, or a disk
does not respond at all, the array therefore \emph{knows which} disk and
which block are affected.  Its redundancy does not have to find the error,
only to reconstruct the contents of a known block.  This is what allows a
single parity block to correct errors in an array (\cref{sec:l17-raid4}),
whereas a parity bit in memory can only detect them.

The RAID variants are called levels and numbered; the number identifies an
organization, not a ranking.  This lesson covers levels 0, 1, 4, 5 and 6.
\margin{In \cite{patterson1988} the \enquote{I} stood for
\emph{inexpensive}.  Levels 2 and 3 are rarely used and are omitted here.}

\begin{keypoint}
  A parity bit only detects an error at an unknown position; a parity block
  reconstructs lost data at a known position.  Disks report their own
  errors, which is what makes parity-based RAID possible.
\end{keypoint}

\section{RAID 0: striping}
\label{sec:l17-raid0}

The first organization uses no redundancy at all and aims only at
performance.\source{Lesson 17, videos 16--18; H\&P §D.2.}  The data of one
logical disk is distributed over several physical disks so that consecutive
blocks are stored on different disks, which is called
\term{striping}[ストライピング].  The blocks at the same position on all
disks form a \emph{stripe}.  An array of $N$ disks that uses striping without
any redundancy is \term{RAID 0}[RAID 0]: logical block $i$ is stored on disk
$i \bmod N$ (\cref{fig:l17-raid01}a).

\begin{figure}[htbp]
  \centering
  % Block placement in (a) RAID 0 with three disks and (b) RAID 1.
% Numbers are logical block numbers.
\begin{tikzpicture}[x=1cm,y=1cm, font=\footnotesize\sffamily]
  \def\blkw{0.95}\def\blkh{0.42}\def\blkgap{1.1}
  % ---------------- (a) RAID 0
  \node[anchor=west, font=\footnotesize\bfseries] at (-0.1,0.95)
    {\iflangja{(a) RAID 0（ストライピング）}{(a) RAID 0 (striping)}};
  \foreach \d in {0,1,2} {
    \node[font=\scriptsize, text=ln-muted] at ({\d*\blkgap+\blkw/2},0.3) {\iflangja{ディスク \d}{Disk \d}};
    \foreach \s in {0,...,3} {
      \pgfmathtruncatemacro{\b}{3*\s+\d}
      \draw[fill=white] ({\d*\blkgap},{-\s*\blkh}) rectangle ++(\blkw,-\blkh);
      \node at ({\d*\blkgap+\blkw/2},{-\s*\blkh-\blkh/2}) {\b};
    }
    \node[font=\scriptsize, text=ln-muted] at ({\d*\blkgap+\blkw/2},{-4*\blkh-0.15}) {$\vdots$};
  }
  % ---------------- (b) RAID 1
  \begin{scope}[xshift=5.6cm]
  \node[anchor=west, font=\footnotesize\bfseries] at (-0.1,0.95)
    {\iflangja{(b) RAID 1（ミラーリング）}{(b) RAID 1 (mirroring)}};
  \foreach \d in {0,1} {
    \node[font=\scriptsize, text=ln-muted] at ({\d*\blkgap+\blkw/2},0.3) {\iflangja{ディスク \d}{Disk \d}};
    \foreach \s in {0,...,3} {
      \draw[fill=white] ({\d*\blkgap},{-\s*\blkh}) rectangle ++(\blkw,-\blkh);
      \node at ({\d*\blkgap+\blkw/2},{-\s*\blkh-\blkh/2}) {\s};
    }
    \node[font=\scriptsize, text=ln-muted] at ({\d*\blkgap+\blkw/2},{-4*\blkh-0.15}) {$\vdots$};
  }
  \end{scope}
\end{tikzpicture}

  \caption{Placement of logical blocks (a) in RAID 0 with three disks, where
    each stripe holds three different blocks, and (b) in RAID 1, where both
    disks hold every block.}
  \label{fig:l17-raid01}
\end{figure}

\subsection{Performance}

The array has $N$ times the capacity of one disk.  Requests for blocks on
different disks are served simultaneously, and a large transfer is split
into $N$ transfers that proceed in parallel, so the throughput is up to $N$
times that of one disk.  An access to a single block still needs a seek and
a rotational latency on its disk (Lesson 16), so small requests complete
sooner only because they spend less time waiting in a queue behind other
requests.  Large transfers also finish sooner, since their data moves to or
from $N$ disks at once.

\subsection{Reliability}

For reliability calculations we assume that disks fail independently of each
other and that each disk has a constant \emph{failure rate}
$f_1 = 1/\mathrm{MTTF}_1$, where $\mathrm{MTTF}_1$ is the MTTF of one
disk.\margin{Real disks fail more often early and late in their life; the
constant rate models the long middle part of their service life.}  The
failure rate of a set of disks is then the sum of the individual rates.  In
RAID 0 each disk holds every $N$-th block of the logical disk, so the failure
of any disk loses data.  The array fails at the rate $N f_1$, and
\begin{equation}
  \mathrm{MTTF}_{\text{RAID 0}} \;=\; \frac{\mathrm{MTTF}_1}{N} .
  \label{eq:l17-raid0}
\end{equation}
For an array, a failure means that data is lost, so its MTTF is also called
the \emph{mean time to data loss} (MTTDL).

\begin{example}
  \label{ex:l17-raid0}
  Disks with $\mathrm{MTTF}_1 = \SI{120000}{\hour}$ (about 13.7 years) are
  combined into a RAID 0 array of five disks.  The array is up to five times
  as fast and as large as one disk, but its MTTF is only
  $\SI{120000}{\hour}/5 = \SI{24000}{\hour}$, about 2.7 years.
\end{example}

\section{RAID 1: mirroring}
\label{sec:l17-raid1}

The second organization spends disks on reliability.\source{Lesson 17,
videos 19--22; H\&P §D.2.}  Writing the same data to two disks, so that each
holds a complete copy of every block, is called
\term{mirroring}[ミラーリング].  An array of two mirrored disks
(\cref{fig:l17-raid01}b) keeps all data available as long as one of its
disks works; it is called \term{RAID 1}[RAID 1].

\begin{itemize}
  \item \textbf{Writes} go to both disks.  The two writes proceed in parallel,
    so writing is as fast as on a single disk.
  \item \textbf{Reads} can be served by either copy, so the read throughput
    is twice that of one disk.
  \item \textbf{Faults} confined to one disk, including the failure of the
    whole disk, are tolerated.  With only two copies a vote is impossible,
    but none is needed: the disk's own code reports which copy is bad, and
    the other copy is used.
  \item \textbf{Cost}: two disks provide the capacity of one.
\end{itemize}

\subsection{Reliability without repair}

While both disks work, the array suffers a disk failure at the rate $2f_1$,
so the first failure comes after $\mathrm{MTTF}_1/2$ on average.  The array
survives it with one disk left.  Because the failure rate is constant, the
surviving disk is as good as new, and it fails after another
$\mathrm{MTTF}_1$ on average.  Only then is data lost:
\begin{equation}
  \mathrm{MTTF}_{\text{RAID 1}} \;=\; \frac{\mathrm{MTTF}_1}{2} + \mathrm{MTTF}_1
  \;=\; 1.5\,\mathrm{MTTF}_1 .
  \label{eq:l17-raid1-norepair}
\end{equation}

\subsection{Reliability with repair}

In practice a failed disk is replaced and the data is copied onto the new
disk from the surviving one.  Let $\mathrm{MTTR}_1$ be the average time from
the failure until the new disk holds a complete copy.  Data is lost only if
the second disk fails within this window (\cref{fig:l17-raid-repair} with
$N=2$).  Since $\mathrm{MTTR}_1 \ll \mathrm{MTTF}_1$, the probability of that
is about $\mathrm{MTTR}_1/\mathrm{MTTF}_1$, so on average
$\mathrm{MTTF}_1/\mathrm{MTTR}_1$ first failures occur before one of them
ends in data loss.  Each is preceded by $\mathrm{MTTF}_1/2$ of both disks
working, hence
\begin{equation}
  \mathrm{MTTF}_{\text{RAID 1}} \;\approx\;
  \frac{\mathrm{MTTF}_1}{2} \times \frac{\mathrm{MTTF}_1}{\mathrm{MTTR}_1}
  \;=\; \frac{\mathrm{MTTF}_1^2}{2\,\mathrm{MTTR}_1} .
  \label{eq:l17-raid1-repair}
\end{equation}

\begin{figure}[htbp]
  \centering
  % States of an array that tolerates one disk failure (RAID 1 with N = 2,
% RAID 4/5 with N disks) when failed disks are replaced.
\begin{tikzpicture}[x=1cm,y=1cm, font=\footnotesize\sffamily, >={Stealth[length=1.8mm]},
  st/.style={draw, rounded corners=3pt, minimum height=0.9cm, text width=2.2cm,
    align=center, inner sep=3pt},
  lab/.style={font=\scriptsize, align=center}]
  \node[st, fill=ln-accent!15] (ok) at (0,0) {\iflangja{$N$ 台すべて動作}{all $N$ disks\\working}};
  \node[st, fill=ln-box] (rep) at (5.4,0) {\iflangja{1 台が故障\\（修復中）}{one disk failed\\(under repair)}};
  \node[st, fill=ln-caution!15] (lost) at (5.4,-2.6) {\iflangja{データ喪失}{data lost}};
  \draw[->] (ok) to[bend left=18] node[lab, above]
    {\iflangja{$N$ 台のいずれかが故障\\（平均 $\mathrm{MTTF}_1/N$ 後）}{one of the $N$ disks fails\\(after $\mathrm{MTTF}_1/N$ on average)}} (rep);
  \draw[->] (rep) to[bend left=18] node[lab, below]
    {\iflangja{修復完了\\（$\mathrm{MTTR}_1$ 後）}{repair completes\\(after $\mathrm{MTTR}_1$)}} (ok);
  \draw[->, ln-caution] (rep) -- node[lab, right, text=black]
    {\iflangja{修復完了前に残り $N-1$ 台の\\いずれかが故障（確率\\$\approx (N-1)\,\mathrm{MTTR}_1/\mathrm{MTTF}_1$）}{one of the other $N-1$ disks\\fails before the repair completes\\(probability $\approx (N-1)\,\mathrm{MTTR}_1/\mathrm{MTTF}_1$)}} (lost);
\end{tikzpicture}

  \caption{An array that tolerates one failed disk, with repair.  Data is
    lost only if a second disk fails while the first is being repaired.  For
    RAID 1, $N=2$.}
  \label{fig:l17-raid-repair}
\end{figure}

\begin{example}
  \label{ex:l17-raid1}
  With the disks of \cref{ex:l17-raid0} ($\mathrm{MTTF}_1 = \SI{120000}{\hour}$),
  a RAID 1 array without repair has an MTTF of \SI{180000}{\hour}, about
  20.5 years.  If a failed disk is replaced and its copy restored within
  $\mathrm{MTTR}_1 = \SI{30}{\hour}$, the MTTF becomes
  $\SI{60000}{\hour} \times 4000 = \SI{2.4e8}{\hour}$, about
  \num{27000} years.
\end{example}

\begin{keypoint}
  Without repair, redundancy helps little: two mirrored disks last only
  $1.5\,\mathrm{MTTF}_1$ on average, because the first of them fails after
  $\mathrm{MTTF}_1/2$.  With
  repair, it multiplies the MTTF by a factor of the order of
  $\mathrm{MTTF}_1/\mathrm{MTTR}_1$.  In a fault-tolerant system, faster
  repair of a module improves the reliability of the system, not only its
  availability.
\end{keypoint}

\section{RAID 4: block-interleaved parity}
\label{sec:l17-raid4}

Parity provides the redundancy of mirroring at a fraction of its
cost.\source{Lesson 17, videos 23--27; H\&P §D.2.}  An array of $N$ disks
can use $N-1$ of them as data disks, striped as in RAID 0, and the remaining
one as a parity disk, whose block in each stripe is the exclusive OR of the
data blocks of that stripe (\cref{fig:l17-raid4-write}a).  This organization
is called \term{RAID 4}[RAID 4].  For a stripe with data blocks
$D_0,\dots,D_{N-2}$, applying \cref{eq:l17-parity} bit by bit gives
\begin{equation}
  P \;=\; D_0 \oplus D_1 \oplus \dots \oplus D_{N-2} .
  \label{eq:l17-raid4-parity}
\end{equation}

Because $x \oplus x = 0$, XORing both sides of \cref{eq:l17-raid4-parity}
with all surviving blocks shows that a lost block is the XOR of all other
blocks of its stripe, parity included.  If a disk fails, every block on it
can be reconstructed from the other $N-1$ disks.  If the failed disk is the
parity disk, the parity is simply recomputed.  The redundancy costs one disk
out of $N$, much less than the second copy of mirroring.

\subsection{Writes}

A write must update the parity of its stripe.  Recomputing it from all data
blocks of the stripe would require reading $N-2$ other disks.  Instead, the
new parity is obtained from the old parity by removing the old data and
adding the new data (\cref{fig:l17-raid4-write}b):
\begin{equation}
  P' \;=\; P \oplus D \oplus D' .
  \label{eq:l17-parity-update}
\end{equation}
A write therefore reads the old data block and the old parity block, and
then writes the new data block and the new parity block: two reads and two
writes, independently of $N$.

\begin{figure}[htbp]
  \centering
  % (a) RAID 4 with three data disks and a parity disk; rows are stripes A-D.
% (b) A small write to block B1: two reads, an XOR, two writes.
\begin{tikzpicture}[x=1cm,y=1cm, font=\footnotesize\sffamily, >={Stealth[length=1.8mm]},
  blk/.style={draw, fill=white, minimum width=1.25cm, minimum height=0.5cm, inner sep=1pt},
  lab/.style={font=\scriptsize, text=ln-muted, align=center}]
  \def\blkw{0.8}\def\blkh{0.45}\def\blkgap{0.9}
  % ---------------- (a) layout
  \node[anchor=west, font=\footnotesize\bfseries] at (-0.9,0.95) {(a) \iflangja{配置}{Layout}};
  \node[font=\scriptsize, text=ln-muted, anchor=east] at (-0.05,0.3) {\iflangja{ディスク}{Disk}};
  \foreach \d in {0,1,2,3} {
    \node[font=\scriptsize, text=ln-muted] at ({\d*\blkgap+\blkw/2},0.3) {\d};
  }
  \foreach \S [count=\r from 0] in {A,B,C,D} {
    \node[font=\scriptsize, text=ln-muted, anchor=east] at (-0.05,{-\r*\blkh-\blkh/2}) {\S};
    \foreach \d in {0,1,2} {
      \draw[fill=white] ({\d*\blkgap},{-\r*\blkh}) rectangle ++(\blkw,-\blkh);
      \node at ({\d*\blkgap+\blkw/2},{-\r*\blkh-\blkh/2}) {$\S_\d$};
    }
    \draw[fill=ln-accent!15] ({3*\blkgap},{-\r*\blkh}) rectangle ++(\blkw,-\blkh);
    \node at ({3*\blkgap+\blkw/2},{-\r*\blkh-\blkh/2}) {$P_\S$};
  }
  \draw[ln-caution, very thick] ({1*\blkgap},{-\blkh}) rectangle ++(\blkw,-\blkh);
  \draw[ln-caution, very thick] ({3*\blkgap},{-\blkh}) rectangle ++(\blkw,-\blkh);
  \node[lab] at ({1*\blkgap+\blkw/2},{-4*\blkh-0.3}) {\iflangja{データ}{data}};
  \node[lab] at ({3*\blkgap+\blkw/2},{-4*\blkh-0.3}) {\iflangja{パリティ}{parity}};
  % ---------------- (b) small write
  \begin{scope}[xshift=5.4cm]
  \node[anchor=west, font=\footnotesize\bfseries] at (-0.4,0.95) {(b) \iflangja{$B_1$ への書き込み}{Write to $B_1$}};
  \node[blk] (nd) at (1.6,0.1)  {\iflangja{新 $B_1'$}{new $B_1'$}};
  \node[blk] (od) at (1.6,-0.7) {\iflangja{旧 $B_1$}{old $B_1$}};
  \node[blk] (op) at (1.6,-1.5) {\iflangja{旧 $P_B$}{old $P_B$}};
  \node[lab, anchor=east] at (0.9,0.1)  {\iflangja{ディスク 1 に\\書く}{written\\to disk 1}};
  \node[lab, anchor=east] at (0.9,-0.7) {\iflangja{ディスク 1 から\\読む}{read from\\disk 1}};
  \node[lab, anchor=east] at (0.9,-1.5) {\iflangja{ディスク 3 から\\読む}{read from\\disk 3}};
  \node[draw, circle, fill=ln-accent!15, inner sep=0pt, minimum size=0.55cm] (x) at (3.3,-0.7) {$\oplus$};
  \draw[->] (nd.east) -- (x);
  \draw[->] (od.east) -- (x);
  \draw[->] (op.east) -- (x);
  \node[blk, fill=ln-accent!15] (np) at (5.1,-0.7) {\iflangja{新 $P_B'$}{new $P_B'$}};
  \draw[->] (x) -- (np);
  \node[lab, anchor=north] at (5.1,-1.0) {\iflangja{ディスク 3 に\\書く}{written\\to disk 3}};
  \end{scope}
\end{tikzpicture}

  \caption{RAID 4.  (a) Three data disks and a parity disk; rows are stripes.
    (b) A write to block $B_1$ reads $B_1$ and $P_B$, computes the new
    parity with \cref{eq:l17-parity-update}, and writes $B_1'$ and $P_B'$.}
  \label{fig:l17-raid4-write}
\end{figure}

\begin{example}
  \label{ex:l17-parity}
  Consider an array with three data disks and a parity disk, and one stripe
  with 4-bit blocks (\cref{tab:l17-parity}).  The parity is
  $1001 \oplus 0111 \oplus 1010 = 0100$.  When disk~1 fails, its block is
  rebuilt as $1001 \oplus 1010 \oplus 0100 = 0111$.  When $0010$ is written
  to disk~1, \cref{eq:l17-parity-update} gives the new parity
  $0100 \oplus 0111 \oplus 0010 = 0001$, which equals the XOR of the new data
  blocks $1001 \oplus 0010 \oplus 1010$.
\end{example}

\begin{table}[htbp]
  \centering\small
  \caption{One stripe of RAID 4 in \cref{ex:l17-parity}.  The reconstructed
    and updated blocks are shown in bold.}
  \label{tab:l17-parity}
  \begin{tabular}{@{}lcccc@{}}
    \toprule
    & Disk 0 & Disk 1 & Disk 2 & Disk 3 (parity) \\
    \midrule
    Initial stripe            & 1001 & 0111 & 1010 & 0100 \\
    Disk 1 lost, rebuilt      & 1001 & \textbf{0111} & 1010 & 0100 \\
    After writing 0010 to disk 1 & 1001 & \textbf{0010} & 1010 & \textbf{0001} \\
    \bottomrule
  \end{tabular}
\end{table}

\subsection{Performance and reliability}

\begin{itemize}
  \item \textbf{Reads} use only the data disks, so the read throughput is
    $N-1$ times that of one disk.
  \item \textbf{Writes} all need the parity disk, once to read and once to
    write.  Writes cannot proceed in parallel even when they go to different
    data disks, and each write occupies the parity disk for two accesses.
    The write throughput is therefore half that of a single disk, whatever
    the value of $N$.
  \item \textbf{Reliability}: the array tolerates the failure of any one
    disk.  Without repair the first failure comes after
    $\mathrm{MTTF}_1/N$ on average and the second after a further
    $\mathrm{MTTF}_1/(N-1)$.  For $N \ge 3$ the sum is less than
    $\mathrm{MTTF}_1$: without repair, such an array is less reliable than a
    single disk.  With repair, the argument of \cref{eq:l17-raid1-repair}
    applies with $N$ disks that can fail first and $N-1$ that can fail
    during the repair (\cref{fig:l17-raid-repair}):
    \begin{equation}
      \mathrm{MTTF}_{\text{RAID 4}} \;\approx\;
      \frac{\mathrm{MTTF}_1}{N} \times
      \frac{\mathrm{MTTF}_1}{(N-1)\,\mathrm{MTTR}_1} .
      \label{eq:l17-raid4}
    \end{equation}
\end{itemize}

\section{RAID 5: distributed parity}
\label{sec:l17-raid5}

The parity disk of RAID 4 is a bottleneck for writes.\source{Lesson 17,
videos 28--29; H\&P §D.2.}  The remedy is to spread the parity over all
disks: each stripe still has one parity block, computed exactly as in RAID 4,
but the parity blocks of different stripes are stored on different disks,
rotating through all $N$ disks (\cref{fig:l17-raid5-6}a).  This organization
is called \term{RAID 5}[RAID 5]; every disk holds both data and parity.

\begin{figure}[htbp]
  \centering
  % (a) RAID 5 with five disks: the parity block rotates among the disks; two
% writes to different stripes use disjoint disks.  (b) RAID 6 with two check
% blocks (P and Q) per stripe.
\begin{tikzpicture}[x=1cm,y=1cm, font=\footnotesize\sffamily]
  \def\blkw{0.86}\def\blkh{0.45}\def\blkgap{0.94}
  % ---------------- (a) RAID 5
  \node[anchor=west, font=\footnotesize\bfseries] at (-0.5,0.95) {(a) RAID 5};
  \node[font=\scriptsize, text=ln-muted, anchor=east] at (-0.05,0.3) {\iflangja{ディスク}{Disk}};
  \foreach \d in {0,...,4} {
    \node[font=\scriptsize, text=ln-muted] at ({\d*\blkgap+\blkw/2},0.3) {\d};
  }
  % stripe letter / parity disk
  \foreach \S/\p [count=\r from 0] in {A/4,B/3,C/2,D/1,E/0} {
    \node[font=\scriptsize, text=ln-muted, anchor=east] at (-0.05,{-\r*\blkh-\blkh/2}) {\S};
    \foreach \d in {0,...,4} {
      \pgfmathtruncatemacro{\i}{\d<\p ? \d : \d-1}
      \ifnum\d=\p
        \draw[fill=ln-accent!15] ({\d*\blkgap},{-\r*\blkh}) rectangle ++(\blkw,-\blkh);
        \node at ({\d*\blkgap+\blkw/2},{-\r*\blkh-\blkh/2}) {$P_\S$};
      \else
        \draw[fill=white] ({\d*\blkgap},{-\r*\blkh}) rectangle ++(\blkw,-\blkh);
        \node at ({\d*\blkgap+\blkw/2},{-\r*\blkh-\blkh/2}) {$\S_\i$};
      \fi
    }
  }
  % write to A0 (disks 0, 4) and to C1 (disks 1, 2)
  \draw[ln-caution, very thick] (0,0) rectangle ++(\blkw,-\blkh);
  \draw[ln-caution, very thick] ({4*\blkgap},0) rectangle ++(\blkw,-\blkh);
  \draw[ln-tip, very thick] ({1*\blkgap},{-2*\blkh}) rectangle ++(\blkw,-\blkh);
  \draw[ln-tip, very thick] ({2*\blkgap},{-2*\blkh}) rectangle ++(\blkw,-\blkh);
  % ---------------- (b) RAID 6
  \begin{scope}[xshift=6.3cm]
  \node[anchor=west, font=\footnotesize\bfseries] at (-0.5,0.95) {(b) RAID 6};
  \node[font=\scriptsize, text=ln-muted, anchor=east] at (-0.05,0.3) {\iflangja{ディスク}{Disk}};
  \foreach \d in {0,...,4} {
    \node[font=\scriptsize, text=ln-muted] at ({\d*\blkgap+\blkw/2},0.3) {\d};
  }
  % stripe letter / P disk / Q disk
  \foreach \S/\p/\q [count=\r from 0] in {A/3/4,B/2/3,C/1/2,D/0/1,E/4/0} {
    \node[font=\scriptsize, text=ln-muted, anchor=east] at (-0.05,{-\r*\blkh-\blkh/2}) {\S};
    \foreach \d in {0,...,4} {
      \pgfmathtruncatemacro{\i}{(\d>\p ? 1 : 0) + (\d>\q ? 1 : 0)}
      \pgfmathtruncatemacro{\i}{\d-\i}
      \ifnum\d=\p
        \draw[fill=ln-accent!15] ({\d*\blkgap},{-\r*\blkh}) rectangle ++(\blkw,-\blkh);
        \node at ({\d*\blkgap+\blkw/2},{-\r*\blkh-\blkh/2}) {$P_\S$};
      \else\ifnum\d=\q
        \draw[fill=ln-accent!30] ({\d*\blkgap},{-\r*\blkh}) rectangle ++(\blkw,-\blkh);
        \node at ({\d*\blkgap+\blkw/2},{-\r*\blkh-\blkh/2}) {$Q_\S$};
      \else
        \draw[fill=white] ({\d*\blkgap},{-\r*\blkh}) rectangle ++(\blkw,-\blkh);
        \node at ({\d*\blkgap+\blkw/2},{-\r*\blkh-\blkh/2}) {$\S_\i$};
      \fi\fi
    }
  }
  \end{scope}
\end{tikzpicture}

  \caption{(a) RAID 5 with five disks.  The parity block rotates among the
    disks; a write to $A_0$ (red) and a write to $C_1$ (green) use disjoint
    disks and can proceed in parallel.  (b) RAID 6 with two check blocks,
    $P$ and $Q$, per stripe.}
  \label{fig:l17-raid5-6}
\end{figure}

\begin{itemize}
  \item \textbf{Reads}: data is spread over all $N$ disks, so the read
    throughput is $N$ times that of one disk.
  \item \textbf{Writes} still need two reads and two writes
    (\cref{eq:l17-parity-update}), but the disk holding the parity depends
    on the stripe.  Writes to different stripes usually involve different
    disks and proceed in parallel.  The four accesses per write are spread
    over all $N$ disks, so the write throughput is $N/4$ times that of one
    disk.
  \item \textbf{Reliability} and \textbf{cost} are the same as for RAID 4:
    one failed disk is tolerated, \cref{eq:l17-raid4} gives the MTTF with
    repair, and one disk's worth of capacity holds parity.
\end{itemize}

\begin{example}
  \label{ex:l17-raid5}
  Six disks, each able to perform 200 accesses per second, with
  $\mathrm{MTTF}_1 = \SI{120000}{\hour}$ and $\mathrm{MTTR}_1 =
  \SI{30}{\hour}$.  As RAID 4 they serve 1000 reads or 100 writes per second;
  as RAID 5, 1200 reads or 300 writes per second.  In both cases the MTTF
  with repair is $(\num{120000}/6) \times \num{120000}/(5 \times 30) =
  \SI{1.6e7}{\hour}$, about \num{1800} years.
\end{example}

\begin{keypoint}
  RAID 5 has the reliability and the cost of RAID 4 and better read and write
  performance.  Since it is at least as good in every respect, RAID 4 is
  rarely used.
\end{keypoint}

\section{RAID 6: two check blocks}
\label{sec:l17-raid6}

RAID 5 tolerates one failed disk; the next organization tolerates
two.\source{Lesson 17, videos 30--32; H\&P §D.2.}  It adds a second check
block to every stripe (\cref{fig:l17-raid5-6}b): one check block, $P$, is the
parity; the other, $Q$, is computed with a different code.  This
organization is called \term{RAID 6}[RAID 6].  $P$ and $Q$ give two
independent equations per stripe, so any two lost blocks of a stripe can be
reconstructed: a single lost block is rebuilt from the parity as in RAID 5,
and two lost blocks are obtained by solving both equations.

\begin{itemize}
  \item \textbf{Reliability}: the array tolerates the failure of any two
    disks.
  \item \textbf{Cost}: two disks' worth of capacity hold check blocks.
  \item \textbf{Writes} must update both check blocks: three reads (old
    data, $P$, $Q$) and three writes, compared with two of each in RAID 5.
    With six accesses per write spread over $N$ disks, the write throughput
    is about $N/6$ times that of one disk; reads remain $N$ times as fast.
\end{itemize}

\subsection{RAID 6 and correlated failures}

For independent disk failures, RAID 5 with repair already achieves an MTTF
of well over a thousand years (\cref{ex:l17-raid5}).  In RAID 6 data is lost
only if a third disk fails while two are being repaired, which multiplies the
MTTF by another factor of the order of
$\mathrm{MTTF}_1/((N-2)\,\mathrm{MTTR}_1)$, about a thousand with the numbers
of the example.  Against independent
failures, RAID 6 is overkill.

Its value lies in failures that are \emph{not} independent of the first
one.  The typical case is an operation fault during the repair: the
operator replaces a working disk instead of the failed one.  A RAID 5 array
now misses two disks and has lost data; a RAID 6 array keeps working and can
still be repaired.

\begin{table}[htbp]
  \centering\small
  \caption{Comparison of the RAID levels for $N$ disks, each with
    $\mathrm{MTTF}_1$ and $\mathrm{MTTR}_1$.  Capacity and throughput are
    relative to one disk.  The MTTF assumes repair for RAID 1, 4, 5 and 6;
    RAID 0 loses data at the first failure.}
  \label{tab:l17-raid}
  \setlength{\tabcolsep}{4pt}
  \begin{tabular}{@{}lccccc@{}}
    \toprule
    Level & Capacity & Reads & Writes & Failures tolerated & MTTF \\
    \midrule
    RAID 0 & $N$   & $N$   & $N$   & 0 & $\mathrm{MTTF}_1/N$ \\
    RAID 1 ($N=2$) & 1 & 2 & 1   & 1 & $\mathrm{MTTF}_1^2/(2\,\mathrm{MTTR}_1)$ \\
    RAID 4 & $N-1$ & $N-1$ & $1/2$ & 1 & $\mathrm{MTTF}_1^2/(N(N-1)\,\mathrm{MTTR}_1)$ \\
    RAID 5 & $N-1$ & $N$   & $N/4$ & 1 & $\mathrm{MTTF}_1^2/(N(N-1)\,\mathrm{MTTR}_1)$ \\
    RAID 6 & $N-2$ & $N$   & $N/6$ & 2 & much higher than RAID 5 \\
    \bottomrule
  \end{tabular}
\end{table}

\begin{keypoint}[Lesson summary]
  A fault in a module becomes an error when activated and a failure when the
  error reaches the delivered service.  Reliability is measured by the MTTF,
  availability by $\mathrm{MTTF}/(\mathrm{MTTF}+\mathrm{MTTR})$.  Faults are
  classified by cause (hardware, design, operation, environment) and by
  duration (permanent, intermittent, transient), and are dealt with by
  avoidance, tolerance and faster repair.  Checkpointing and $N$-module
  redundancy tolerate faults in general; parity and ECC protect memory; RAID
  protects against failed disks.  Striping (RAID 0) trades reliability for
  speed, mirroring (RAID 1) buys reliability with a second copy, and parity
  (RAID 4, 5) does so at the cost of one disk, with RAID 5 avoiding the
  parity-disk bottleneck.  With repair, the MTTF of a redundant array grows
  with $\mathrm{MTTF}_1/\mathrm{MTTR}_1$ (\cref{tab:l17-raid}); RAID 6 guards
  mainly against correlated failures such as mistakes during repair.
\end{keypoint}

\end{document}
